\section{Purpose and Scope}

This specification describes the C implementation of \code{fix-hostfiles} as
an algebra of state transformations.  It is not a transcription of the source
code.  Implementation details appear only where they refine, constrain, or
qualify the abstract behavior.

The program coordinates four operations around \code{hblock(1)} and the macOS
hosts file:

\begin{description}
  \item[prep] Preserve the active hosts file and run \code{hblock}.
  \item[restore] Replace the active hosts file with the preserved original.
  \item[add] Add a DNS name to the allow list and remove its active mappings.
  \item[flush] Flush the DNS cache and restart \code{mDNSResponder}.
\end{description}

The algebra has four purposes:

\begin{enumerate}
  \item define the state visible at the program boundary;
  \item state precise laws for exact matching, idempotence, commutativity, and
        reversibility;
  \item distinguish the laws of primitive transformations from the weaker
        guarantees of multi-file commands; and
  \item provide properties that can be checked independently of individual
        implementation examples.
\end{enumerate}

The specification deliberately separates several kinds of equality.  An
operation can be idempotent with respect to the active hosts file and allow
list while still changing backup history.  Similarly, an atomic rename can
protect one pathname without making the complete command transactional.

\subsection{Notation}

Function application is written as $f(x)$.  Composition is written as
$f \circ g$, so $(f \circ g)(x)=f(g(x))$.  A finite sequence is written using
angle brackets, and $\concat$ denotes concatenation.  State field replacement
uses the notation

\[
  S[H \coloneqq H', A \coloneqq A'].
\]

Every field not named in a replacement retains its old value.  Equations are
interpreted over successful executions unless a failure result is shown
explicitly.
